%
% FH Technikum Wien
% !TEX encoding = UTF-8 Unicode
%
% Version V2.24 von 2024-12-19 otrebski
%
\documentclass[BMR,Seminar,ngerman,IEEE]{twbook}
\renewcommand*{\degreecourse}{FH Technikum Wien, Bachelorstudiengang BIF\\Lehrveranstaltung BIF-VZ-5-WS2025-INFC-EN}

% --- Encoding & fonts ---
\usepackage[utf8]{inputenc}
\usepackage[T1]{fontenc}

% --- Graphics, figures, urls, lists, tables, code ---
\usepackage{float}
\usepackage{placeins}
\usepackage{graphicx}
\usepackage{caption}
\usepackage{booktabs}
\usepackage{tabularx}
\usepackage{enumitem}
\usepackage{listings}
\usepackage[final]{microtype} % hilft gegen Overfull/Underfull
\emergencystretch=2em         % sanfter Puffer für Zeilenumbruch

% Screenshots hier hinein legen:
\graphicspath{{screenshots/}}

% --- Bibliography (twbook lädt biblatex via Hook, sobald inputenc geladen ist) ---
\addbibresource{Literatur.bib}

% --- Listings (inkl. YAML) ---
\lstdefinelanguage{yaml}{
  keywords={true,false,null,y,n},
  comment=[l]{\#},
  morestring=[b]',
  morestring=[b]",
  sensitive=false,
  showstringspaces=false,
  literate =    {---}{{\textcolor{gray}{---}}}3
                {>}{{\textcolor{gray}{>}}}1
                {|}{{\textcolor{gray}{|}}}1
                {:}{{\textcolor{blue}{:}}}1
                {-}{{\textcolor{gray}{-}}}1,
}
\lstset{
  basicstyle=\ttfamily\small,
  frame=single,
  breaklines=true,
  backgroundcolor=\color{gray!5},
  captionpos=b
}

% --- Convenience figure macros ---
\newcommand{\screenshotH}[3]{%
  \begin{figure}[H]
    \centering
    \includegraphics[width=\linewidth]{#1}%
    \caption{#2}%
    \label{fig:#3}%
  \end{figure}%
}
\newcommand{\placeholderfig}[2]{%
  \begin{figure}[htbp]
    \centering
    \fbox{\rule{0pt}{120pt}\rule{0.9\linewidth}{0pt}}%
    \caption{#1}%
    \label{fig:#2}%
  \end{figure}%
}

% Die nachfolgenden Pakete stellen sonst nicht benötigte Features zur Verfügung
\usepackage{blindtext}

% --- Hyperref am Schluss laden (robuster) ---
\usepackage{hyperref}

%
% Einträge für Deckblatt, Kurzfassung, etc.
%
\title{Terraform Workshop Protocol}
\author{David Veigel \and Kevin Forter}
\studentnumber{XXXXXXXXXXXXXXX} % ggf. beide Nummern: 12345678, 87654321
\place{Wien}
% \acknowledgements{\blindtext}

\begin{document}

\maketitle

% ============================================================================
\onecolumn
\newpage

% ============================================================================
\chapter{Workshop 1: Terraform Basics}
In this workshop, we set up a simple web server on AWS using Terraform. The goal is to provision an EC2 instance in the default VPC, install an Apache web server using user data, and verify that it is reachable.

\section{Configuration}

\subsection{Provider and Region}
We configure the AWS provider to use the \texttt{us-east-1} region, as specified in the requirements.
\begin{lstlisting}[language=bash]
provider "aws" {
  region = "us-east-1"
}
\end{lstlisting}

\subsection{VPC and Subnets}
Instead of creating a new network, we use the existing default VPC and its subnets. This simplifies the setup for this initial workshop.
\begin{lstlisting}[language=bash]
data "aws_vpc" "default" {
  default = true
}

data "aws_subnets" "default_vpc_subnets" {
  filter {
    name   = "vpc-id"
    values = [data.aws_vpc.default.id]
  }
}
\end{lstlisting}

\subsection{Security Group}
To make our web server accessible, we create a security group. We allow inbound TCP traffic on port 80 (HTTP) from any IP address (\texttt{0.0.0.0/0}). We also allow all outbound traffic to ensure the server can download updates and packages.
\begin{lstlisting}[language=bash]
resource "aws_security_group" "web_sg" {
  name        = "tf-web-sg"
  description = "Allow HTTP in, all egress"
  vpc_id      = data.aws_vpc.default.id

  ingress {
    description      = "HTTP"
    from_port        = 80
    to_port          = 80
    protocol         = "tcp"
    cidr_blocks      = ["0.0.0.0/0"]
  }

  egress {
    description      = "All egress"
    from_port        = 0
    to_port          = 0
    protocol         = "-1"
    cidr_blocks      = ["0.0.0.0/0"]
  }
}
\end{lstlisting}
\screenshotH{ws_1/terraform_ws_1_security_group.png}{Security Group Configuration in AWS Console}{sg-ws1}

\subsection{EC2 Instance}
We provision an EC2 instance using the \texttt{t2.micro} instance type, which is eligible for the free tier. We dynamically fetch the latest Ubuntu 22.04 LTS AMI. The \texttt{user\_data} script is used to automatically install and start the Apache web server upon launch.
\begin{lstlisting}[language=bash]
resource "aws_instance" "web" {
  ami                         = data.aws_ami.ubuntu.id
  instance_type               = "t2.micro"
  subnet_id                   = data.aws_subnets.default_vpc_subnets.ids[0]
  vpc_security_group_ids      = [aws_security_group.web_sg.id]
  associate_public_ip_address = true

  user_data = <<-EOF
    #!/bin/bash
    sudo apt-get update
    sudo apt-get install -y apache2
    sudo systemctl start apache2
    sudo systemctl enable apache2
    echo "<h1>Hello World</h1>" | sudo tee /var/www/html/index.html
  EOF

  tags = {
    Name = "tf-apache"
  }
}
\end{lstlisting}
\screenshotH{ws_1/terraform_ws_1_instances.png}{EC2 Instance Running in AWS Console}{ec2-ws1}

\subsection{Output and Verification}
Finally, we output the public DNS of the instance to easily access it. We also include a \texttt{null\_resource} to verify that the web server is actually reachable via HTTP.
\begin{lstlisting}[language=bash]
output "instance_public_dns" {
  description = "Public DNS of the EC2 instance"
  value       = aws_instance.web.public_dns
}
\end{lstlisting}

\section{Execution}
We initialize the Terraform project to download the necessary providers.
\screenshotH{ws_1/terraform_ws_1_init.png}{Terraform Init}{init-ws1}

Then we apply the configuration to create the resources.
\screenshotH{ws_1/terraform_ws_1_auto_approve.png}{Terraform Apply}{apply-ws1}

% ============================================================================
\chapter{Workshop 2: Advanced Networking}
In the second workshop, we move away from the default VPC and build a custom network infrastructure. This includes creating a VPC, subnets, routing tables, and an Internet Gateway. We also assign a persistent Elastic IP to our web server.

\section{Network Infrastructure}

\subsection{VPC and Internet Gateway}
We create a new VPC with the CIDR block \texttt{10.0.0.0/16}. To enable communication with the internet, we attach an Internet Gateway to this VPC.
\begin{lstlisting}[language=bash]
resource "aws_vpc" "web_vpc" {
  cidr_block           = "10.0.0.0/16"
  enable_dns_support   = true
  enable_dns_hostnames = true
  tags = { Name = "tf-web-vpc" }
}

resource "aws_internet_gateway" "igw" {
  vpc_id = aws_vpc.web_vpc.id
  tags = { Name = "tf-web-igw" }
}
\end{lstlisting}
\screenshotH{ws_2/vpc.png}{Custom VPC Created}{vpc-ws2}
\screenshotH{ws_2/internet_gateway.png}{Internet Gateway Attached}{igw-ws2}

\subsection{Subnet and Routing}
We create a subnet (\texttt{10.0.1.0/24}) within our VPC. To allow traffic from the subnet to reach the internet, we create a custom route table that routes all traffic (\texttt{0.0.0.0/0}) to the Internet Gateway, and associate it with our subnet.
\begin{lstlisting}[language=bash]
resource "aws_route_table" "web_rt" {
  vpc_id = aws_vpc.web_vpc.id
  route {
    cidr_block = "0.0.0.0/0"
    gateway_id = aws_internet_gateway.igw.id
  }
  tags = { Name = "tf-web-rt" }
}

resource "aws_subnet" "web_subnet" {
  vpc_id                  = aws_vpc.web_vpc.id
  cidr_block              = "10.0.1.0/24"
  availability_zone       = "us-east-1a"
  map_public_ip_on_launch = false
  tags = { Name = "tf-web-subnet" }
}

resource "aws_route_table_association" "web_rta" {
  subnet_id      = aws_subnet.web_subnet.id
  route_table_id = aws_route_table.web_rt.id
}
\end{lstlisting}
\screenshotH{ws_2/subnet_routing_table.png}{Subnet and Route Table Association}{rt-ws2}

\section{Compute and Access}

\subsection{Security Group}
Similar to Workshop 1, we create a security group to allow HTTP traffic.
\screenshotH{ws_2/security_group_ingres.png}{Security Group Inbound Rules}{sg-ws2}

\subsection{Network Interface and Elastic IP}
Instead of relying on auto-assigned public IPs, we create a Network Interface (ENI) with a fixed private IP (\texttt{10.0.1.50}) and associate an Elastic IP (EIP) with it. This gives our server a static public IP address.
\begin{lstlisting}[language=bash]
resource "aws_network_interface" "web_eni" {
  subnet_id       = aws_subnet.web_subnet.id
  private_ips     = ["10.0.1.50"]
  security_groups = [aws_security_group.web_sg.id]
  tags = { Name = "tf-web-eni" }
}

resource "aws_eip" "web_eip" {
  domain                    = "vpc"
  network_interface         = aws_network_interface.web_eni.id
  associate_with_private_ip = "10.0.1.50"
  depends_on                = [aws_internet_gateway.igw]
  tags = { Name = "tf-web-eip" }
}
\end{lstlisting}
\screenshotH{ws_2/elastic_ip.png}{Elastic IP Allocation}{eip-ws2}

\subsection{EC2 Instance}
We launch the EC2 instance and attach the previously created network interface. This ensures the instance uses our specific network configuration.
\begin{lstlisting}[language=bash]
resource "aws_instance" "web" {
  ami           = data.aws_ami.ubuntu.id
  instance_type = "t2.micro"

  network_interface {
    network_interface_id = aws_network_interface.web_eni.id
    device_index         = 0
  }

  user_data = <<-EOF
    #!/bin/bash
    sudo apt-get update
    sudo apt-get install -y apache2
    sudo systemctl start apache2
    sudo systemctl enable apache2
    echo "<h1>Hello World</h1>" | sudo tee /var/www/html/index.html
  EOF

  tags = { Name = "tf-apache-instance" }
}
\end{lstlisting}
\screenshotH{ws_2/ec2_instance_summary.png}{EC2 Instance Summary}{ec2-ws2}

\section{Execution}
We apply the configuration and verify the creation of resources. The output confirms the public IP of our instance.
\screenshotH{ws_2/terraform_apply_1.png}{Terraform Apply Execution}{apply-ws2}

\end{document}
